
\documentclass[11pt,a4paper]{article}

\usepackage[margin=1in]{geometry}
\usepackage[T1]{fontenc}
\usepackage[utf8]{inputenc}
\usepackage{lmodern}
\usepackage{microtype}
\usepackage{parskip}
\usepackage{enumitem}
\usepackage{longtable}
\usepackage{booktabs}
\usepackage{array}
\usepackage{xcolor}
\usepackage{hyperref}
\usepackage{listings}
\usepackage{titlesec}
\usepackage{fancyhdr}

\definecolor{darkblue}{RGB}{25,55,95}
\definecolor{lightgray}{RGB}{245,245,245}
\definecolor{codegray}{RGB}{90,90,90}

\hypersetup{
    colorlinks=true,
    linkcolor=darkblue,
    urlcolor=darkblue,
    pdftitle={MediFlow ERP - Complete Project Task Description},
    pdfauthor={MediFlow ERP Project}
}

\pagestyle{fancy}
\fancyhf{}
\lhead{MediFlow ERP}
\rhead{Project Task Description}
\cfoot{\thepage}

\titleformat{\section}{\Large\bfseries\color{darkblue}}{\thesection}{0.6em}{}
\titleformat{\subsection}{\large\bfseries\color{darkblue}}{\thesubsection}{0.6em}{}
\titleformat{\subsubsection}{\normalsize\bfseries}{\thesubsubsection}{0.6em}{}

\setlist[itemize]{leftmargin=*, itemsep=2pt, topsep=3pt}
\setlist[enumerate]{leftmargin=*, itemsep=3pt, topsep=3pt}

\lstset{
    backgroundcolor=\color{lightgray},
    basicstyle=\ttfamily\small,
    frame=single,
    breaklines=true,
    columns=fullflexible
}

\newcommand{\task}[1]{\item[$\square$] #1}
\newcommand{\important}[1]{\textbf{Important:} #1}

\begin{document}

\begin{titlepage}
    \centering
    \vspace*{2cm}
    {\Huge\bfseries MediFlow ERP\par}
    \vspace{0.5cm}
    {\LARGE Complete Project Task Description\par}
    \vspace{1.5cm}
    {\large Enterprise Healthcare ERP for Private Hospitals and Clinics\par}
    \vfill

    \begin{minipage}{0.9\textwidth}
    \textbf{Project objective.}
    MediFlow is an enterprise-grade ERP designed to unify clinical,
    logistical, scheduling, and financial operations for private hospital
    networks and clinics. It acts as a single source of truth across
    patient care, pharmacy and supply chain, scheduling, billing,
    insurance, and accounting.
    \end{minipage}

    \vfill
    {\large Development Roadmap and Implementation Checklist\par}
\end{titlepage}

\tableofcontents
\newpage

% ============================================================
\section{Project Overview}

MediFlow addresses a common healthcare operational problem: data is often
siloed between reception, clinical staff, pharmacy, inventory, billing,
insurance, and accounting. This can cause delays, inconsistent information,
stock allocation errors, financial leakage, and weak accountability.

The system must unify these domains while maintaining strict access control,
transactional consistency, detailed auditability, and strong protection of
sensitive patient and financial information.

\subsection{Primary Business Domains}

The four core business domains are:

\begin{enumerate}
    \item \textbf{Patient Care and Admissions} --- \texttt{apps.clinical}
    \item \textbf{Pharmacy and Supply Chain} --- \texttt{apps.inventory}
    \item \textbf{Rostering and Resource Allocation} --- \texttt{apps.scheduling}
    \item \textbf{Billing and Insurance Ledger} --- \texttt{apps.finance}
\end{enumerate}

Supporting applications should also be introduced:

\begin{itemize}
    \item \texttt{apps.accounts} --- authentication, users, roles, permissions
    \item \texttt{apps.organizations} --- organizations, hospitals, departments
    \item \texttt{apps.audit} --- immutable audit history
    \item \texttt{apps.notifications} --- notifications and background workflows
    \item \texttt{apps.core} --- shared infrastructure and common functionality
\end{itemize}

\subsection{Core Cross-Module Workflow}

The principal ERP workflow is:

\begin{lstlisting}
Patient Checks In
        |
        v
Doctor Consultation
        |
        v
Doctor Prescribes Medication
        |
        v
Inventory Batch Selection
        |
        v
Database Concurrency Lock
        |
        v
Stock Reservation / Deduction
        |
        v
Medication Dispensed
        |
        v
Invoice Item Generated
        |
        v
Accounting Ledger Updated
        |
        v
Audit Log Created
\end{lstlisting}

A critical requirement is that related operations are performed atomically
where business consistency requires it.

% ============================================================
\section{Phase 0 --- Project Definition and Architecture}

\subsection{0.1 Define Project Scope}

\task{Define the MVP scope.}
\task{Define features planned for post-MVP releases.}
\task{Define the hospital and clinic organizational structure.}
\task{Define supported user types.}
\task{Define the permissions of every user type.}
\task{Define which modules each role can access.}
\task{Document all major business workflows.}
\task{Identify regulatory, privacy, accounting, and operational requirements.}

\subsection{0.2 Define User Types}

At minimum, define:

\begin{itemize}
    \item Super Administrator
    \item Hospital Administrator
    \item Doctor
    \item Nurse
    \item Receptionist
    \item Pharmacist
    \item Accountant
    \item Insurance Manager
\end{itemize}

The exact role hierarchy should be documented before implementation.

\subsection{0.3 Define Core Domain Entities}

Create the initial domain map:

\begin{lstlisting}
Organization
  |
  +-- Hospital
  |     +-- Department
  |     +-- Ward
  |     +-- Room
  |     +-- Operating Room
  |
  +-- Users
  +-- Patients
  +-- Appointments
  +-- Admissions
  +-- Prescriptions
  +-- Inventory
  +-- Finance
\end{lstlisting}

\subsection{0.4 Define Application Boundaries}

The initial Django structure should be:

\begin{lstlisting}
apps/
    accounts/
    organizations/
    clinical/
    inventory/
    scheduling/
    finance/
    audit/
    notifications/
    core/
\end{lstlisting}

Define which model and business responsibility belongs to each app.
Avoid creating circular dependencies between applications.

\subsection{0.5 Define Cross-Module Dependencies}

Document the dependency direction:

\begin{lstlisting}
accounts
    |
    v
organizations
    |
    +----> scheduling
    |
    +----> clinical
                |
                v
            inventory
                |
                v
             finance
                |
                v
              audit
\end{lstlisting}

\subsection{0.6 Define API Architecture}

\task{Define REST API conventions.}
\task{Define API versioning.}
\task{Define authentication strategy.}
\task{Define pagination conventions.}
\task{Define filtering conventions.}
\task{Define searching conventions.}
\task{Define ordering conventions.}
\task{Define standard error responses.}
\task{Define serializer validation strategy.}
\task{Define OpenAPI documentation strategy.}

% ============================================================
\section{Phase 1 --- Infrastructure and Django Foundation}

\subsection{1.1 Repository}

\task{Create the Git repository.}
\task{Create a production-safe \texttt{.gitignore}.}
\task{Create \texttt{.env.example}.}
\task{Create the project README.}
\task{Document local development instructions.}
\task{Establish a branching strategy.}

Recommended branches:

\begin{lstlisting}
main
develop
feature/*
bugfix/*
\end{lstlisting}

\subsection{1.2 Django Project}

\task{Create the Django project.}
\task{Configure Django settings.}
\task{Separate base, development, production, and testing settings.}
\task{Configure installed applications.}
\task{Configure middleware.}
\task{Configure templates and static files.}
\task{Configure media files.}
\task{Configure timezone.}
\task{Configure internationalization where required.}

\subsection{1.3 Core Dependencies}

Evaluate and install the required packages, including:

\begin{lstlisting}
Django
Django REST Framework
django-filter
psycopg
drf-spectacular
django-simple-history
Celery
Redis
pytest
pytest-django
factory-boy
\end{lstlisting}

Potential supporting packages:

\begin{lstlisting}
django-cors-headers
django-health-check
gunicorn
\end{lstlisting}

Every dependency must have a documented reason for inclusion.

\subsection{1.4 PostgreSQL}

\task{Configure PostgreSQL.}
\task{Create the development database.}
\task{Configure Django's database connection.}
\task{Configure connection pooling if required.}
\task{Verify transaction behavior.}
\task{Run initial migrations.}
\task{Verify rollback behavior.}

\subsection{1.5 Docker Development Environment}

Create a development environment containing:

\begin{lstlisting}
Django
PostgreSQL
Redis
Celery Worker
Celery Beat
Optional Mailpit / development mail server
\end{lstlisting}

The target is:

\begin{lstlisting}
docker compose up
\end{lstlisting}

to start the complete local development environment.

% ============================================================
\section{Phase 2 --- Organization and User Management}

This phase must be completed before most clinical, inventory, and finance
functionality because these domains depend on users, roles, and
organizational context.

\subsection{2.1 Custom User}

Create the custom \texttt{User} model with appropriate fields, for example:

\begin{lstlisting}
id
email
password
first_name
last_name
phone
is_active
is_staff
created_at
updated_at
\end{lstlisting}

\task{Use the custom user model from the beginning of the project.}
\task{Make email uniqueness behavior explicit.}
\task{Implement secure password hashing.}
\task{Implement account activation/deactivation.}

\subsection{2.2 RBAC Models}

Create:

\begin{lstlisting}
Role
Permission
UserRole
\end{lstlisting}

Define permissions such as:

\begin{lstlisting}
clinical.read
clinical.create_prescription
clinical.update_prescription
patient.read
inventory.read
inventory.receive
inventory.dispense
finance.read
finance.create_invoice
finance.process_payment
\end{lstlisting}

\subsection{2.3 Organizations}

Create:

\begin{lstlisting}
Organization
Hospital
Department
\end{lstlisting}

Relationships:

\begin{lstlisting}
Organization
    |
    v
Hospital
    |
    v
Department
\end{lstlisting}

\subsection{2.4 Staff Assignment}

Create:

\begin{lstlisting}
StaffProfile
DoctorProfile
NurseProfile
\end{lstlisting}

\task{Associate staff with organizations.}
\task{Associate staff with hospitals.}
\task{Associate staff with departments.}
\task{Define employment/active status.}

\subsection{2.5 Authentication}

Implement:

\task{Login.}
\task{Logout.}
\task{JWT or another explicitly selected authentication mechanism.}
\task{Access-token expiration.}
\task{Refresh tokens if JWT is selected.}
\task{Password reset.}
\task{Email verification.}
\task{Account activation and deactivation.}

\subsection{2.6 DRF Permission System}

Implement reusable permission classes such as:

\begin{lstlisting}
IsDoctor
IsPharmacist
IsAccountant
IsHospitalAdmin
\end{lstlisting}

Also implement object-level and organization-level authorization.

% ============================================================
\section{Phase 3 --- Audit System}

The audit system should be implemented before sensitive clinical and
financial workflows.

\subsection{3.1 Audit Model}

Create:

\begin{lstlisting}
AuditLog
\end{lstlisting}

Track:

\begin{lstlisting}
user
action
model
object_id
timestamp
IP address
user agent
old_values
new_values
\end{lstlisting}

\subsection{3.2 Audit Events}

Support at least:

\begin{lstlisting}
CREATE
UPDATE
DELETE
VIEW
LOGIN
LOGOUT
PRESCRIBE
DISPENSE
PAYMENT
CLAIM
\end{lstlisting}

\subsection{3.3 Immutable Audit Trail}

\task{Prevent ordinary users from editing audit records.}
\task{Prevent ordinary users from deleting audit records.}
\task{Restrict audit-log access to authorized roles.}
\task{Record the acting user.}
\task{Record timestamps.}
\task{Record affected objects.}
\task{Record old values.}
\task{Record new values.}
\task{Document retention and archival strategy.}

Use \texttt{django-simple-history}, custom audit infrastructure, or another
documented approach. Database-level protections/triggers may be considered
for stronger immutability requirements.

\subsection{3.4 Audit Tests}

Verify:

\begin{lstlisting}
Doctor changes prescription
        |
        v
AuditLog
        |
        +-- Who?
        +-- When?
        +-- What changed?
        +-- Old value?
        +-- New value?
\end{lstlisting}

% ============================================================
\section{Phase 4 --- Patient Management}

Implement inside \texttt{apps.clinical}.

\subsection{4.1 Patient Model}

Create \texttt{Patient}.

Potential fields:

\begin{lstlisting}
medical_record_number
first_name
last_name
date_of_birth
sex
phone
email
address
emergency_contact
blood_type
status
created_at
updated_at
\end{lstlisting}

\subsection{4.2 Patient Identification}

\task{Generate a unique Medical Record Number (MRN).}
\task{Search by MRN.}
\task{Search by name.}
\task{Search by phone number.}
\task{Implement duplicate-patient detection.}
\task{Document patient identity matching rules.}

\subsection{4.3 Patient Demographics}

Implement:

\task{Personal information.}
\task{Contact information.}
\task{Emergency contact.}
\task{Allergies.}
\task{Medical history.}

\subsection{4.4 Patient Privacy}

Explicitly define which patient fields are visible to:

\begin{itemize}
    \item Doctor
    \item Nurse
    \item Receptionist
    \item Pharmacist
    \item Accountant
    \item Administrator
\end{itemize}

Implement queryset-level and field-level access controls where necessary.

% ============================================================
\section{Phase 5 --- Scheduling}

Implement \texttt{apps.scheduling}.

\subsection{5.1 Doctor Availability}

Create:

\begin{lstlisting}
DoctorSchedule
WorkingHours
TimeOff
\end{lstlisting}

\task{Define working hours.}
\task{Define recurring schedules if required.}
\task{Define leave/time-off.}
\task{Define exceptional working days.}

\subsection{5.2 Appointment}

Create:

\begin{lstlisting}
Appointment
\end{lstlisting}

Fields should include:

\begin{lstlisting}
patient
doctor
department
start_time
end_time
status
reason
notes
\end{lstlisting}

Statuses:

\begin{lstlisting}
SCHEDULED
CONFIRMED
CHECKED_IN
IN_PROGRESS
COMPLETED
CANCELLED
NO_SHOW
\end{lstlisting}

\subsection{5.3 Appointment Operations}

\task{Create appointment.}
\task{Reschedule appointment.}
\task{Cancel appointment.}
\task{Confirm appointment.}
\task{Check patient in.}
\task{Mark appointment as in progress.}
\task{Complete appointment.}
\task{Record no-show.}

\subsection{5.4 Appointment Conflict Prevention}

Prevent:

\begin{lstlisting}
Doctor
    +-- 10:00 Patient A
    +-- 10:00 Patient B  <-- invalid
\end{lstlisting}

Use database constraints and/or transactional service logic.

\subsection{5.5 Operating Rooms}

Create:

\begin{lstlisting}
OperatingRoom
OperatingRoomSchedule
\end{lstlisting}

\task{Create operating rooms.}
\task{Track availability.}
\task{Reserve rooms.}
\task{Prevent overlapping reservations.}
\task{Support cancellation and rescheduling.}

% ============================================================
\section{Phase 6 --- Patient Check-In and Admissions}

\subsection{6.1 Check-In Workflow}

Implement:

\begin{lstlisting}
Appointment
    |
    v
Patient arrives
    |
    v
Check-in
    |
    v
Waiting
    |
    v
Doctor consultation
\end{lstlisting}

\subsection{6.2 Admission}

Create:

\begin{lstlisting}
Admission
Ward
Room
Bed
\end{lstlisting}

Relationship:

\begin{lstlisting}
Patient
    |
    v
Admission
    |
    v
Ward
    |
    v
Room
    |
    v
Bed
\end{lstlisting}

\subsection{6.3 Bed Management}

Support statuses:

\begin{lstlisting}
AVAILABLE
OCCUPIED
RESERVED
MAINTENANCE
\end{lstlisting}

\task{Assign bed.}
\task{Release bed.}
\task{Transfer patient.}
\task{Prevent double assignment.}
\task{Track bed history.}

\subsection{6.4 Discharge}

Implement:

\begin{lstlisting}
Admission
    |
    v
Discharge
    |
    +-- Final clinical summary
    |
    +-- Final billing
\end{lstlisting}

% ============================================================
\section{Phase 7 --- EHR and Clinical Records}

\subsection{7.1 Encounter}

Create:

\begin{lstlisting}
Encounter
\end{lstlisting}

An encounter represents a clinical interaction between a patient and
clinical staff.

\begin{lstlisting}
Patient
    |
    v
Encounter
    |
    v
Doctor
\end{lstlisting}

\subsection{7.2 Clinical Notes}

Create:

\begin{lstlisting}
ClinicalNote
\end{lstlisting}

Support:

\begin{itemize}
    \item Chief complaint
    \item History
    \item Examination
    \item Assessment
    \item Plan
\end{itemize}

\subsection{7.3 Diagnosis}

Create:

\begin{lstlisting}
Diagnosis
\end{lstlisting}

Design the model so that a standardized diagnosis coding system such as
ICD-10 can be added or integrated later.

\subsection{7.4 Vitals}

Create:

\begin{lstlisting}
VitalSign
\end{lstlisting}

Track:

\begin{lstlisting}
temperature
heart_rate
blood_pressure
respiratory_rate
oxygen_saturation
weight
height
\end{lstlisting}

\task{Record who measured the vital sign.}
\task{Record measurement time.}
\task{Maintain historical measurements.}

\subsection{7.5 Medical History}

Support:

\task{Chronic conditions.}
\task{Previous surgeries.}
\task{Allergies.}
\task{Previous admissions.}
\task{Family history.}
\task{Other clinically relevant history.}

% ============================================================
\section{Phase 8 --- Medication Catalog and Prescriptions}

\subsection{8.1 Medication Catalog}

Create:

\begin{lstlisting}
Medication
MedicationForm
DosageUnit
\end{lstlisting}

Support different forms, for example:

\begin{lstlisting}
Tablet
Syrup
Injection
Capsule
Drops
Cream
Spray
Solution
Powder
Suppository
\end{lstlisting}

\subsection{8.2 Prescription}

Create:

\begin{lstlisting}
Prescription
PrescriptionItem
\end{lstlisting}

Structure:

\begin{lstlisting}
Prescription
    |
    +-- Patient
    +-- Doctor
    +-- Encounter
    |
    +-- Items
          +-- Medication
          +-- Dose
          +-- Frequency
          +-- Duration
          +-- Quantity
\end{lstlisting}

\subsection{8.3 Prescription Lifecycle}

Implement:

\begin{lstlisting}
DRAFT
  |
  v
ISSUED
  |
  v
DISPENSING
  |
  v
DISPENSED
\end{lstlisting}

Also support:

\begin{lstlisting}
CANCELLED
EXPIRED
\end{lstlisting}

\subsection{8.4 Prescription Auditing}

\task{Audit creation.}
\task{Audit modification.}
\task{Audit cancellation.}
\task{Record the responsible doctor/user.}
\task{Protect finalized prescription records from unauthorized modification.}

% ============================================================
\section{Phase 9 --- Inventory and Pharmacy}

Implement \texttt{apps.inventory}.

\subsection{9.1 Product Catalog}

Create:

\begin{lstlisting}
Product
ProductCategory
Manufacturer
Supplier
\end{lstlisting}

Support both medication and medical supplies.

\subsection{9.2 Batch Tracking}

Create:

\begin{lstlisting}
InventoryBatch
\end{lstlisting}

Track:

\begin{lstlisting}
product
batch_number
manufacture_date
expiration_date
purchase_price
selling_price
quantity
\end{lstlisting}

Also consider:

\begin{lstlisting}
warehouse
storage_location
unit_of_measure
received_at
status
\end{lstlisting}

\subsection{9.3 Warehouse and Storage}

Create:

\begin{lstlisting}
Warehouse
StorageLocation
\end{lstlisting}

\task{Create warehouses.}
\task{Create storage locations.}
\task{Transfer stock between locations.}
\task{Track current stock by location.}

\subsection{9.4 Stock Ledger}

Do not rely only on a mutable quantity field.

Create:

\begin{lstlisting}
StockMovement
\end{lstlisting}

Movement types:

\begin{lstlisting}
RECEIPT
TRANSFER
ADJUSTMENT
RESERVATION
DISPENSING
RETURN
EXPIRATION
\end{lstlisting}

Every stock change should have a traceable movement record.

% ============================================================
\section{Phase 10 --- Inventory Concurrency and Race Conditions}

This is a critical technical requirement.

\subsection{10.1 Transactional Allocation}

Use Django transactions:

\begin{lstlisting}
transaction.atomic()
\end{lstlisting}

and row-level database locking:

\begin{lstlisting}
select_for_update()
\end{lstlisting}

Conceptual workflow:

\begin{lstlisting}
Prescription
    |
    v
Find eligible batch
    |
    v
BEGIN TRANSACTION
    |
    v
SELECT ... FOR UPDATE
    |
    v
Check quantity
    |
    v
Reserve / deduct quantity
    |
    v
Create StockMovement
    |
    v
COMMIT
\end{lstlisting}

\subsection{10.2 Race-Condition Test}

Simulate:

\begin{lstlisting}
Doctor A ----\
              >---- Last available vial
Doctor B ----/
\end{lstlisting}

Expected behavior:

\begin{lstlisting}
One request succeeds.
The other request waits for the transaction and then
fails safely because the stock is no longer available.
\end{lstlisting}

The implementation must not permit negative or double-allocated stock.

\subsection{10.3 Expiration Handling}

Implement:

\task{FEFO (First Expire, First Out).}
\task{Expired batch detection.}
\task{Blocking expired stock from dispensing.}
\task{Near-expiration warnings.}
\task{Expiration stock movements.}

\subsection{10.4 Temperature-Sensitive Inventory}

Support:

\begin{lstlisting}
TemperatureControlled
StorageCondition
TemperatureLog
\end{lstlisting}

Example storage range:

\begin{lstlisting}
2 C -- 8 C
\end{lstlisting}

\task{Record temperature measurements.}
\task{Detect out-of-range values.}
\task{Generate alerts where required.}
\task{Track affected stock if a cold-chain violation occurs.}

% ============================================================
\section{Phase 11 --- Prescription to Pharmacy Integration}

Connect clinical and inventory.

Workflow:

\begin{lstlisting}
Doctor
    |
    v
Prescription
    |
    v
Pharmacy
    |
    v
Batch selection
    |
    v
Stock reservation
    |
    v
Dispensing
    |
    v
Stock movement
\end{lstlisting}

Tasks:

\task{Allow pharmacists to view eligible prescriptions.}
\task{Find available inventory batches.}
\task{Apply FEFO selection.}
\task{Reserve stock.}
\task{Dispense medication.}
\task{Create stock movement records.}
\task{Update prescription status.}
\task{Audit the dispensing event.}

% ============================================================
\section{Phase 12 --- Finance and Billing}

Implement \texttt{apps.finance}.

\subsection{12.1 Price Catalog}

Create:

\begin{lstlisting}
Service
ServicePrice
MedicationPrice
\end{lstlisting}

Support services such as:

\begin{lstlisting}
Consultation
MRI
Blood Test
Medication
Hospital Bed
Surgery
\end{lstlisting}

\task{Support price changes over time.}
\task{Preserve the price actually used on historical invoices.}

\subsection{12.2 Patient Billing Folio}

Create:

\begin{lstlisting}
PatientFolio
\end{lstlisting}

The folio represents the patient's active financial account for a
clinical encounter or admission.

\subsection{12.3 Invoice}

Create:

\begin{lstlisting}
Invoice
InvoiceItem
\end{lstlisting}

Example:

\begin{lstlisting}
Invoice
    +-- Consultation     $50
    +-- Medication       $20
    +-- Lab test         $30
    +-- Room             $100
\end{lstlisting}

Implement:

\task{Invoice creation.}
\task{Invoice itemization.}
\task{Tax/discount rules if required.}
\task{Invoice status.}
\task{Invoice finalization.}
\task{Invoice cancellation rules.}
\task{Invoice audit trail.}

% ============================================================
\section{Phase 13 --- Clinical to Finance Integration}

Implement automatic billing triggers.

Consultation:

\begin{lstlisting}
Patient
    |
    v
Consultation
    |
    v
InvoiceItem
\end{lstlisting}

Medication:

\begin{lstlisting}
Prescription
    |
    v
Dispensing
    |
    v
InvoiceItem
\end{lstlisting}

The system must ensure that a successful dispensing operation creates the
appropriate financial transaction without creating duplicate charges.

Define idempotency rules for retryable operations.

% ============================================================
\section{Phase 14 --- Insurance}

\subsection{14.1 Insurance Providers}

Create:

\begin{lstlisting}
InsuranceProvider
InsurancePlan
\end{lstlisting}

\subsection{14.2 Patient Insurance}

Create:

\begin{lstlisting}
PatientInsurance
\end{lstlisting}

Track relevant policy information, effective dates, status, and coverage.

\subsection{14.3 Coverage Rules}

Support:

\begin{lstlisting}
Insurance coverage percentage
Patient responsibility
Annual deductible
Coverage limit
Copayment
Excluded services
\end{lstlisting}

Example:

\begin{lstlisting}
80% insurance
20% patient
\end{lstlisting}

\subsection{14.4 Claims}

Create:

\begin{lstlisting}
InsuranceClaim
InsuranceClaimItem
\end{lstlisting}

Workflow:

\begin{lstlisting}
Invoice
    |
    v
Claim
    |
    v
Submitted
    |
    +----> Approved
    |
    +----> Rejected
    |
    v
Payment / Reconciliation
\end{lstlisting}

\task{Create claim.}
\task{Submit claim.}
\task{Track claim status.}
\task{Record rejection reason.}
\task{Support resubmission where applicable.}
\task{Reconcile insurer payments.}

% ============================================================
\section{Phase 15 --- Accounting Ledger}

This phase turns the finance module into an ERP-style accounting system.

\subsection{15.1 Accounting Models}

Create:

\begin{lstlisting}
Account
JournalEntry
JournalEntryLine
\end{lstlisting}

\subsection{15.2 Chart of Accounts}

Example:

\begin{lstlisting}
1000 Assets
1100 Cash
1200 Accounts Receivable
1300 Inventory

2000 Liabilities
2100 Insurance Payable

4000 Revenue
4100 Consultation Revenue
4200 Pharmacy Revenue
\end{lstlisting}

The actual chart of accounts must be defined with the accounting
requirements of the target organization/jurisdiction.

\subsection{15.3 Double-Entry Accounting}

Every financial transaction must satisfy:

\[
\text{Total Debits} = \text{Total Credits}
\]

Example medication transaction:

\begin{lstlisting}
Medication dispensed
    |
    +-- Debit:  Patient Receivable
    |
    +-- Credit: Pharmacy Inventory
\end{lstlisting}

\task{Implement balanced journal entries.}
\task{Prevent unbalanced entries.}
\task{Audit financial entries.}
\task{Define reversal instead of destructive modification for finalized entries.}

% ============================================================
\section{Phase 16 --- Payments}

Create:

\begin{lstlisting}
Payment
PaymentMethod
Refund
\end{lstlisting}

Support:

\begin{lstlisting}
Cash
Card
Bank Transfer
Insurance
\end{lstlisting}

Workflow:

\begin{lstlisting}
Invoice
    |
    v
Payment
    |
    v
Receipt
    |
    v
Ledger Entry
\end{lstlisting}

Tasks:

\task{Record payments.}
\task{Generate receipts.}
\task{Support partial payments if required.}
\task{Support refunds.}
\task{Prevent duplicate payment processing.}
\task{Create corresponding accounting entries.}
\task{Audit all payment operations.}

% ============================================================
\section{Phase 17 --- End-to-End Cross-Module Transaction}

Implement the primary MediFlow transaction.

\begin{lstlisting}
Patient checks in
        |
        v
Doctor consultation
        |
        v
Doctor prescribes medication
        |
        v
Pharmacist receives prescription
        |
        v
Database transaction
        |
        v
SELECT FOR UPDATE
        |
        v
Batch locked
        |
        v
Stock reserved
        |
        v
Medication dispensed
        |
        v
StockMovement created
        |
        v
InvoiceItem created
        |
        v
JournalEntry created
        |
        v
AuditLog created
        |
        v
COMMIT
\end{lstlisting}

This workflow must be implemented through a dedicated service/domain layer,
not through a large DRF view.

% ============================================================
\section{Phase 18 --- Notifications and Background Jobs}

Create \texttt{apps.notifications} and configure Celery.

\subsection{18.1 Notifications}

Implement support for:

\task{Appointment reminders.}
\task{Prescription notifications.}
\task{Low-stock alerts.}
\task{Expiration alerts.}
\task{Insurance claim updates.}
\task{Payment reminders.}

\subsection{18.2 Scheduled Jobs}

Examples:

\begin{lstlisting}
Every morning
    -> check expiring medication

Every hour
    -> check low-stock products

Daily
    -> appointment reminders

Weekly
    -> financial reports
\end{lstlisting}

\task{Configure Celery worker.}
\task{Configure Celery Beat.}
\task{Define retry policy.}
\task{Handle failed jobs.}
\task{Make important jobs idempotent.}
\task{Monitor task failures.}

% ============================================================
\section{Phase 19 --- Reporting and Dashboards}

\subsection{19.1 Clinical Dashboard}

Display:

\begin{lstlisting}
Today's patients
Today's appointments
Admissions
Discharges
Active patients
\end{lstlisting}

\subsection{19.2 Pharmacy Dashboard}

Display:

\begin{lstlisting}
Total inventory
Low stock
Expiring soon
Expired stock
Dispensed today
\end{lstlisting}

\subsection{19.3 Finance Dashboard}

Display:

\begin{lstlisting}
Revenue
Outstanding invoices
Payments
Insurance claims
Rejected claims
Accounts receivable
\end{lstlisting}

\subsection{19.4 Management Dashboard}

Combine major operational indicators:

\begin{lstlisting}
Patients
Appointments
Admissions
Revenue
Inventory value
Low-stock alerts
Insurance activity
\end{lstlisting}

All dashboard data must respect the user's permissions and organizational
scope.

% ============================================================
\section{Phase 20 --- Security and Privacy}

Security must be treated as a continuous requirement rather than a final
checkbox.

\subsection{20.1 Authentication Security}

\task{Enforce strong password policy.}
\task{Use secure password hashing.}
\task{Configure token expiration.}
\task{Implement secure session management.}
\task{Implement login throttling/rate limiting.}
\task{Consider account lockout policies.}

\subsection{20.2 Authorization}

Verify examples such as:

\begin{lstlisting}
Accountant -> invoice       YES
Accountant -> diagnosis     NO

Doctor -> diagnosis         YES
Doctor -> accounting ledger NO

Pharmacist -> inventory     YES
Pharmacist -> diagnosis     NO
\end{lstlisting}

Do not rely solely on frontend hiding. Enforce authorization on the
backend/API.

\subsection{20.3 Field-Level and Object-Level Access}

For example:

\begin{lstlisting}
Patient
    +-- demographics -> receptionist: YES
    +-- diagnosis    -> receptionist: NO
    +-- invoice      -> accountant: YES
    +-- prescription -> pharmacist: YES
\end{lstlisting}

Implement queryset filtering and field-level serialization where necessary.

\subsection{20.4 API Security}

\task{Implement object-level permissions.}
\task{Implement organization-level data isolation.}
\task{Implement rate limiting where required.}
\task{Validate all input.}
\task{Prevent unauthorized object access.}
\task{Use Django/DRF protections against injection attacks.}
\task{Configure CSRF correctly.}
\task{Configure CORS correctly.}
\task{Configure secure HTTP headers.}
\task{Protect secrets using environment variables or a secret manager.}
\task{Never expose sensitive data in logs.}

\subsection{20.5 Healthcare Privacy}

Before production deployment, identify the applicable legal and regulatory
requirements for the target jurisdiction. Do not assume that merely using
Django or encryption makes a system compliant.

Implement appropriate technical controls for:

\begin{itemize}
    \item least-privilege access
    \item auditability
    \item data minimization
    \item encryption in transit
    \item encryption at rest where required
    \item secure backups
    \item access monitoring
    \item retention and deletion policies
\end{itemize}

% ============================================================
\section{Phase 21 --- Testing}

This project requires extensive automated testing.

\subsection{21.1 Unit Tests}

Test:

\task{Models.}
\task{Services.}
\task{Validators.}
\task{Permission classes.}
\task{Business calculations.}
\task{Financial calculations.}
\task{Inventory calculations.}

\subsection{21.2 API Tests}

Test all major:

\begin{lstlisting}
POST
GET
PATCH
DELETE
\end{lstlisting}

operations, including authentication and authorization behavior.

\subsection{21.3 RBAC Tests}

Example:

\begin{lstlisting}
Doctor -> patient data             YES
Doctor -> accounting data          NO
Pharmacist -> inventory            YES
Accountant -> clinical diagnosis   NO
\end{lstlisting}

Create tests for every sensitive role/module combination.

\subsection{21.4 Concurrency Tests}

Simulate:

\begin{lstlisting}
100 concurrent requests
        |
        v
1 available medication unit
\end{lstlisting}

Expected result:

\begin{lstlisting}
Only one successful allocation.
No negative stock.
No duplicate allocation.
All unsuccessful requests receive a safe error.
\end{lstlisting}

\subsection{21.5 Accounting Tests}

Verify:

\[
\text{Debit} = \text{Credit}
\]

for every finalized journal entry.

\subsection{21.6 End-to-End Workflow Tests}

Test:

\begin{lstlisting}
Patient
    |
    v
Appointment
    |
    v
Check-in
    |
    v
Consultation
    |
    v
Prescription
    |
    v
Dispensing
    |
    v
Invoice
    |
    v
Payment
\end{lstlisting}

% ============================================================
\section{Phase 22 --- Performance}

\subsection{22.1 Database Performance}

\task{Analyze slow queries.}
\task{Add appropriate indexes.}
\task{Use \texttt{select\_related} correctly.}
\task{Use \texttt{prefetch\_related} correctly.}
\task{Eliminate N+1 query problems.}
\task{Configure database connection pooling where appropriate.}
\task{Use database constraints to enforce important invariants.}

\subsection{22.2 Caching}

Use Redis where appropriate for:

\begin{lstlisting}
Frequently accessed data
Sessions
Rate limiting
Temporary calculations
\end{lstlisting}

Do not cache sensitive data without a clear invalidation and authorization
strategy.

\subsection{22.3 Load Testing}

Test progressively:

\begin{lstlisting}
100 concurrent users
500 concurrent users
1000 concurrent users
\end{lstlisting}

Pay particular attention to:

\begin{lstlisting}
Inventory allocation
Appointment booking
Billing
Authentication
High-volume reporting
\end{lstlisting}

% ============================================================
\section{Phase 23 --- Observability}

Implement:

\begin{lstlisting}
Application logs
Error tracking
Performance monitoring
Audit logs
Health checks
\end{lstlisting}

Create health endpoints:

\begin{lstlisting}
/health/
/health/db/
/health/redis/
\end{lstlisting}

Monitor:

\begin{lstlisting}
Request latency
5xx errors
Database errors
Celery failures
Queue length
Background job failures
\end{lstlisting}

Use structured logging where practical.

% ============================================================
\section{Phase 24 --- API Documentation}

Use OpenAPI/Swagger through a documented DRF-compatible solution.

Document endpoints such as:

\begin{lstlisting}
/api/v1/auth/

/api/v1/patients/
/api/v1/appointments/
/api/v1/encounters/
/api/v1/prescriptions/

/api/v1/inventory/
/api/v1/batches/
/api/v1/dispensing/

/api/v1/invoices/
/api/v1/payments/
/api/v1/insurance/

/api/v1/reports/
\end{lstlisting}

For each endpoint document:

\task{Authentication requirements.}
\task{Required permissions.}
\task{Request parameters.}
\task{Request body.}
\task{Validation rules.}
\task{Successful response.}
\task{Error responses.}
\task{Pagination/filtering behavior where applicable.}

% ============================================================
\section{Phase 25 --- Deployment}

Recommended production architecture:

\begin{lstlisting}
                    Internet
                       |
                       v
                 Reverse Proxy
                       |
              +--------+--------+
              |                 |
              v                 v
          Django API         Frontend
              |
       +------+------+
       |             |
       v             v
 PostgreSQL        Redis
                     |
                +----+----+
                |         |
                v         v
             Celery    Celery
             Worker     Beat
\end{lstlisting}

Tasks:

\task{Create production settings.}
\task{Configure production environment variables.}
\task{Configure production PostgreSQL.}
\task{Configure production Redis.}
\task{Configure Celery worker.}
\task{Configure Celery Beat.}
\task{Configure HTTPS.}
\task{Configure domain and DNS.}
\task{Configure static files.}
\task{Configure media storage.}
\task{Configure database backups.}
\task{Configure monitoring.}
\task{Configure log collection.}
\task{Run production migrations safely.}
\task{Document rollback procedures.}

% ============================================================
\section{Phase 26 --- Backup and Disaster Recovery}

Because MediFlow stores clinical and financial data, backup and recovery
must be treated as first-class functionality.

\task{Configure automated PostgreSQL backups.}
\task{Define backup retention periods.}
\task{Store backups securely.}
\task{Test restoration from backup.}
\task{Define point-in-time recovery strategy where required.}
\task{Document disaster recovery procedures.}
\task{Document database migration recovery procedures.}
\task{Define Recovery Point Objective (RPO).}
\task{Define Recovery Time Objective (RTO).}

% ============================================================
\section{Phase 27 --- Final End-to-End Validation}

\subsection{27.1 Outpatient Scenario}

Run the complete workflow:

\begin{lstlisting}
Register patient
    |
    v
Book appointment
    |
    v
Check in
    |
    v
Doctor consultation
    |
    v
Diagnosis
    |
    v
Prescription
    |
    v
Dispensing
    |
    v
Invoice
    |
    v
Payment
\end{lstlisting}

Verify that all clinical, inventory, financial, and audit records are
consistent.

\subsection{27.2 Inpatient Scenario}

Run:

\begin{lstlisting}
Patient
    |
    v
Admission
    |
    v
Bed assignment
    |
    v
Daily treatment
    |
    v
Medication
    |
    v
Laboratory procedures
    |
    v
Room charges
    |
    v
Discharge
    |
    v
Final invoice
    |
    v
Insurance claim
\end{lstlisting}

\subsection{27.3 Inventory Race Condition Scenario}

Run:

\begin{lstlisting}
Doctor A ----\
              >---- 1 available vial
Doctor B ----/
\end{lstlisting}

Verify:

\begin{lstlisting}
Exactly one allocation succeeds.
The other transaction fails safely.
Stock never becomes negative.
No duplicate allocation exists.
Audit records are correct.
Financial records are correct.
\end{lstlisting}

\subsection{27.4 Permission Isolation Scenario}

Verify:

\begin{lstlisting}
Accountant
    -> Invoice: YES
    -> Diagnosis: NO

Doctor
    -> Diagnosis: YES
    -> Ledger: NO

Pharmacist
    -> Inventory: YES
    -> Unauthorized clinical information: NO
\end{lstlisting}

% ============================================================
\section{Recommended Development Order}

The following sequence should be used as the main implementation order:

\begin{enumerate}
    \item Project architecture
    \item Django foundation
    \item PostgreSQL and Docker
    \item Custom User
    \item Organizations
    \item RBAC
    \item Authentication
    \item Audit system
    \item Patient management
    \item Scheduling
    \item Appointments
    \item Patient check-in
    \item Admissions
    \item EHR
    \item Vitals
    \item Diagnoses
    \item Medication catalog
    \item Prescriptions
    \item Inventory catalog
    \item Warehouses
    \item Inventory batches
    \item Stock ledger
    \item Concurrency and locking
    \item Prescription-to-inventory integration
    \item Dispensing
    \item Finance catalog
    \item Patient folio
    \item Invoices
    \item Inventory-to-finance integration
    \item Insurance
    \item Payments
    \item Accounting ledger
    \item Full clinical-to-inventory-to-finance integration
    \item Notifications
    \item Celery
    \item Reports
    \item Security hardening
    \item Automated tests
    \item Performance optimization
    \item Observability
    \item API documentation
    \item Deployment
    \item Backups
    \item Disaster recovery
    \item End-to-end testing
    \item Production release
\end{enumerate}

% ============================================================
\section{Business Logic Architecture}

A major architectural rule for MediFlow is:

\begin{center}
\Large
\textbf{Do not place complex business logic directly inside DRF views.}
\end{center}

Use a layered architecture:

\begin{lstlisting}
HTTP Request
    |
    v
DRF View / ViewSet
    |
    v
Serializer
    |
    v
Service / Domain Layer
    |
    v
Models / Repository / Database
\end{lstlisting}

For complex workflows, the service layer should own transaction boundaries
and business invariants.

\subsection{Example: Dispensing Service}

A conceptual implementation should look like:

\begin{lstlisting}
POST /prescriptions/{id}/dispense/
            |
            v
DispensePrescriptionService
            |
            v
transaction.atomic()
            |
            v
select_for_update()
            |
            v
Validate prescription
            |
            v
Find eligible batch
            |
            v
Check available quantity
            |
            v
Reserve / deduct stock
            |
            v
Create StockMovement
            |
            v
Create InvoiceItem
            |
            v
Create JournalEntry
            |
            v
Create AuditLog
            |
            v
COMMIT
\end{lstlisting}

This structure keeps concurrency handling, financial consistency, inventory
logic, and auditing together as one controlled business operation.

% ============================================================
\section{Definition of Done}

A feature should not be considered complete merely because its endpoint
works.

For every significant feature, verify:

\begin{enumerate}
    \item Database model is correct.
    \item Database constraints are defined.
    \item Business rules are implemented in the service/domain layer.
    \item API endpoint is implemented.
    \item Authentication is enforced.
    \item Authorization is enforced.
    \item Input validation is implemented.
    \item Audit behavior is implemented.
    \item Automated tests are written.
    \item Error handling is implemented.
    \item Documentation is updated.
    \item Performance has been considered.
    \item Cross-module effects have been tested.
\end{enumerate}

For sensitive operations, additionally verify:

\begin{itemize}
    \item Transaction boundaries
    \item Concurrency behavior
    \item Idempotency
    \item Auditability
    \item Data visibility
    \item Financial consistency
\end{itemize}

% ============================================================
\section{MediFlow Completion Checklist}

\subsection*{Foundation}
\begin{itemize}
    \task{Repository configured}
    \task{Django project configured}
    \task{PostgreSQL configured}
    \task{Redis configured}
    \task{Docker environment configured}
    \task{Development and production settings separated}
\end{itemize}

\subsection*{Identity and Access}
\begin{itemize}
    \task{Custom User implemented}
    \task{Roles implemented}
    \task{Permissions implemented}
    \task{Organization isolation implemented}
    \task{Authentication implemented}
    \task{Authorization tested}
\end{itemize}

\subsection*{Clinical}
\begin{itemize}
    \task{Patient management}
    \task{Appointments}
    \task{Check-in}
    \task{Admissions}
    \task{Beds}
    \task{Encounters}
    \task{Clinical notes}
    \task{Diagnoses}
    \task{Vitals}
    \task{Medical history}
    \task{Medication catalog}
    \task{Prescriptions}
\end{itemize}

\subsection*{Inventory}
\begin{itemize}
    \task{Products}
    \task{Categories}
    \task{Manufacturers}
    \task{Suppliers}
    \task{Warehouses}
    \task{Storage locations}
    \task{Inventory batches}
    \task{Expiration tracking}
    \task{Stock ledger}
    \task{FEFO}
    \task{Stock reservation}
    \task{Dispensing}
    \task{Concurrency locking}
    \task{Temperature-controlled inventory}
\end{itemize}

\subsection*{Finance}
\begin{itemize}
    \task{Services and prices}
    \task{Patient folio}
    \task{Invoices}
    \task{Invoice items}
    \task{Insurance providers}
    \task{Insurance plans}
    \task{Patient insurance}
    \task{Claims}
    \task{Payments}
    \task{Refunds}
    \task{Chart of accounts}
    \task{Journal entries}
    \task{Double-entry accounting}
\end{itemize}

\subsection*{Cross-Module Integration}
\begin{itemize}
    \task{Clinical to inventory}
    \task{Inventory to finance}
    \task{Clinical to finance}
    \task{Insurance to finance}
    \task{Audit across all modules}
    \task{Notifications}
\end{itemize}

\subsection*{Engineering Quality}
\begin{itemize}
    \task{Unit tests}
    \task{API tests}
    \task{RBAC tests}
    \task{Concurrency tests}
    \task{Accounting tests}
    \task{End-to-end tests}
    \task{Performance tests}
    \task{Logging}
    \task{Monitoring}
    \task{Health checks}
    \task{OpenAPI documentation}
\end{itemize}

\subsection*{Production}
\begin{itemize}
    \task{Production configuration}
    \task{HTTPS}
    \task{Domain}
    \task{Database backups}
    \task{Backup restoration tested}
    \task{Disaster recovery plan}
    \task{Monitoring}
    \task{Deployment documentation}
    \task{Rollback procedure}
    \task{Final security review}
    \task{Final end-to-end validation}
\end{itemize}

% ============================================================
\section{Final Architecture Target}

The completed MediFlow system should conceptually provide:

\begin{lstlisting}
                         MediFlow ERP
                              |
       +----------------------+----------------------+
       |                      |                      |
       v                      v                      v
   Accounts              Organizations           Audit
       |                      |                      |
       +----------------------+----------------------+
                              |
          +-------------------+-------------------+
          |                   |                   |
          v                   v                   v
      Clinical           Scheduling           Inventory
          |                   |                   |
          |                   |                   |
          +-------------------+-------------------+
                              |
                              v
                           Finance
                              |
                    +---------+---------+
                    |                   |
                    v                   v
                Insurance           Accounting
\end{lstlisting}

The system's central design principle is:

\begin{center}
\textbf{
One patient record, one operational source of truth,
transactionally consistent workflows, strict authorization,
and complete accountability.
}
\end{center}

\end{document}
